\documentclass[11pt]{article}
\usepackage[utf8]{inputenc}        % AJOUT : encodage UTF-8
\usepackage{vmargin}
\setpapersize{A4}
\setmargins{15mm}{10mm}{180mm}{252mm}%
             {16mm}{2mm}{0pt}{10mm}
\usepackage[french]{babel}
\usepackage{lmodern}
\usepackage[T1]{fontenc}
\usepackage{fouriernc}
\usepackage{eurosym}
\usepackage{graphicx}
\usepackage{array}
\usepackage{tabularx}
\usepackage{variations}
\usepackage{amsmath}
\usepackage{amssymb}
\usepackage{epic}
\usepackage{rotating}
\usepackage[autolanguage]{numprint}
\usepackage[tikz]{bclogo}
\usepackage{pstricks}
\usepackage{pstricks-add}
\usepackage{tikz}                  % AJOUT : pour tikzpicture
\usepackage{pifont}                % AJOUT : pour \ding
\usepackage{multirow}
\usepackage{multicol}
\usepackage{setspace}
\usepackage{fancyhdr}
\usepackage{graphpap}
\usepackage{lastpage}

\setlength{\unitlength}{1mm}
\usepackage{xy}
\xyoption{all}

\newcommand{\cadreblanc}[1]{%
\medskip\framebox[\linewidth][c]{%
\rule{0mm}{#1mm}}\par
\medskip}

\def\labelenumii{\theenumi.~\theenumii.}
\def\theenumii{\arabic{enumii}}

\presetkeys{bclogo}{logo=\bcquestion,couleur=gray!10,barre=snake,arrondi=0.1}{}

%%%%%%%%%%%%%%%%%%%%%%%%%%%%%%%%%%%%%%%%%%%%%%%%%%%%%%%%%%
% Types de colonnes AVEC argument (pour tabular ordinaire)
%%%%%%%%%%%%%%%%%%%%%%%%%%%%%%%%%%%%%%%%%%%%%%%%%%%%%%%%%%
\newcolumntype{Z}[1]{>{\raggedright\arraybackslash}m{#1}}
\newcolumntype{Y}[1]{>{\centering\arraybackslash}m{#1}}

%%%%%%%%%%%%%%%%%%%%%%%%%%%%%%%%%%%%%%%%%%%%%%%%%%%%%%%%%%
% Types de colonnes SANS argument (pour tabularx)
%%%%%%%%%%%%%%%%%%%%%%%%%%%%%%%%%%%%%%%%%%%%%%%%%%%%%%%%%%
\newcolumntype{C}{>{\centering\arraybackslash}X}
\newcolumntype{L}{>{\raggedright\arraybackslash}X}
\newcolumntype{R}{>{\raggedleft\arraybackslash}X}

\renewcommand{\thesection}{{\arabic{section}.}}
\addto\captionsfrench{\def\tablename{Document}}

\renewcommand{\footrulewidth}{0.4pt}

\newcommand{\appel}{\raisebox{-2ex}{\includegraphics[width=1cm]{images/appel.jpg}}\quad}

%%%%%%%%%%%%%%%%%%%%%%%%%%%%%%%%%%%%%%%%%%%%%%%%%%%%%%%%%%
% Évaluation par compétences (marginpar)
%%%%%%%%%%%%%%%%%%%%%%%%%%%%%%%%%%%%%%%%%%%%%%%%%%%%%%%%%%
\newcommand\SA{\marginpar{\scalebox{.8}{\framebox[1.6ex]{\rule{0ex}{1.6ex}\raisebox{.75ex}{\rule{1.6ex}{.5pt}}}\hspace{-.3pt}\framebox[3.2ex]{S}}}}
\newcommand\A{\marginpar{\scalebox{.8}{\framebox[1.6ex]{\rule{0ex}{1.6ex}\raisebox{.75ex}{\rule{1.6ex}{.5pt}}}\hspace{-.3pt}\framebox[3.2ex]{A}}}}
\newcommand\Rcomp{\marginpar{\scalebox{.8}{\framebox[1.6ex]{\rule{0ex}{1.6ex}\raisebox{.75ex}{\rule{1.6ex}{.5pt}}}\hspace{-.3pt}\framebox[3.2ex]{R}}}}  % renommé Rcomp
\newcommand\V{\marginpar{\scalebox{.8}{\framebox[1.6ex]{\rule{0ex}{1.6ex}\raisebox{.75ex}{\rule{1.6ex}{.5pt}}}\hspace{-.3pt}\framebox[3.2ex]{V}}}}
\newcommand\Ccomp{\marginpar{\scalebox{.8}{\framebox[1.6ex]{\rule{0ex}{1.6ex}\raisebox{.75ex}{\rule{1.6ex}{.5pt}}}\hspace{-.3pt}\framebox[3.2ex]{C}}}}  % renommé Ccomp

%%%%%%%%%%%%%%%%%%%%%%%%%%%%%%%%%%%%%%%%%%%%%%%%%%%%%%%%%%
% Définitions année, spécialité, diplôme...
%%%%%%%%%%%%%%%%%%%%%%%%%%%%%%%%%%%%%%%%%%%%%%%%%%%%%%%%%%
\newcommand{\specialite}{MICROTECHNIQUES}
\newcommand{\annee}{2025|26}
\newcommand{\diplome}{Baccalauréat Professionnel}
\newcommand{\etablissement}{LP Denis Diderot}
\newcommand{\classe}{TMICRO}
\newcommand{\datedeval}{13 mai 2026}
\newcommand{\matiere}{Sciences Physiques et chimiques}
\newcommand{\sequence}{2/2}

%%%%%%%%%%%%%%%%%%%%%%%%%%%%%%%%%%%%%%%%%%%%%%%%%%%%%%%%%%
% Capacités, connaissances, attitudes
%%%%%%%%%%%%%%%%%%%%%%%%%%%%%%%%%%%%%%%%%%%%%%%%%%%%%%%%%%
\newcommand{\theme}{Caractériser la propagation d'un signal sonore.}

\newcommand{\capacites}{
  \begin{minipage}[c]{12.5cm}
    \begin{itemize}
      \item \rule{0ex}{2.5ex}Exploiter la relation liant la vitesse de propagation, la longueur d'onde et la fréquence d'une onde sonore.
      \item Déterminer expérimentalement la vitesse de propagation d'un son dans l'air ou dans l'eau.\rule[-1.5ex]{0ex}{3ex}
    \end{itemize}
  \end{minipage}
}

\newcommand{\connaissances}{
  \begin{minipage}[c]{12.5cm}
    \begin{itemize}
      \item \rule{0ex}{2.5ex}Savoir que la propagation d'un son nécessite un milieu matériel.
      \item Savoir que la vitesse du son dépend du milieu de propagation.
      \item Connaître la relation qui lie la longueur d'onde, la vitesse de propagation et la période d'une onde sonore $\lambda = c \times T$.
      \item Connaître les ordres de grandeur des vitesses de propagation du son dans l'air et dans l'eau.\rule[-1.5ex]{0ex}{3ex}
    \end{itemize}
  \end{minipage}
}

\newcommand{\attitudes}{
  \begin{minipage}[c]{12.5cm}
    \begin{itemize}
      \item \rule{0ex}{2.5ex}Le sens de l'observation.
      \item L'ouverture à la communication, au dialogue et au débat argumenté.
      \item La rigueur et la précision.
      \item Le respect des règles élémentaires de sécurité.
      \item \rule[-1.5ex]{0ex}{3ex}Le goût de chercher et de raisonner.
    \end{itemize}
  \end{minipage}
}

%%%%%%%%%%%%%%%%%%%%%%%%%%%%%%%%%%%%%%%%%%%%%%%%%%%%%%%%%%
%                   Début du document
%%%%%%%%%%%%%%%%%%%%%%%%%%%%%%%%%%%%%%%%%%%%%%%%%%%%%%%%%%

\begin{document}
\pagestyle{fancy}
\lfoot{Alexis Ruiz}
\cfoot{\thepage/4}
\rfoot{\datedeval}

%%%%%%%%%%%%%%%%%%%%%%%%%%%%%%%%%%%%%%%%%%%%%%%%%%%%%%%%%%
%                     Page de garde
%%%%%%%%%%%%%%%%%%%%%%%%%%%%%%%%%%%%%%%%%%%%%%%%%%%%%%%%%%

\thispagestyle{empty}
\Large
\noindent\begin{tabularx}{1\linewidth}[c]{|c|c|c|C|C|}
  \hline
  \multirow{2}{*}{
  \begin{minipage}[c]{3cm}
    \includegraphics[width=30mm]{images/aca-besancon.png}
  \end{minipage}}
  &
  \multicolumn{4}{c|}{
  \begin{minipage}[c]{12cm}
    \begin{center}
      \rule[-1.5ex]{0ex}{5ex}BACCALAURÉAT PROFESSIONNEL\\
      \rule[-2ex]{0ex}{4ex}\specialite
    \end{center}
  \end{minipage}}\\
  \cline{2-5}
  & \multicolumn{3}{c|}{\rule[-2ex]{0ex}{5.5ex}Épreuve E1 -- Sciences Physiques et chimiques} & Coef.~1.5\\
  \hline
  \multirow{2}{*}{
  \begin{minipage}[c]{3cm}
    \begin{center}
      Contrôle en \\ Cours de\\ Formation
    \end{center}
  \end{minipage}}
  &
  \multirow{2}{*}{
  \raisebox{-2ex}{\begin{minipage}[c]{6.3cm}
    \begin{center}
      Situation d'évaluation de\\ \matiere
    \end{center}
  \end{minipage}}}
  &
  \multirow{2}{*}{
  \begin{minipage}[c]{2.2cm}
    \begin{center}
      Année\\ scolaire\\ \annee
    \end{center}
  \end{minipage}}
  &
  \rule[-1ex]{0ex}{4ex}Séquence &
  Durée\\
  \cline{4-5}
  &&&
  \rule[-1.5ex]{0ex}{5ex}\sequence
  &
  45~min\\
  \hline
\end{tabularx}

\bigskip

\noindent\fbox{
  \begin{minipage}[c]{.97\linewidth}
    \centerline{\rule[-.8ex]{0ex}{3.5ex}FICHE D'INFORMATION DU CANDIDAT}
  \end{minipage}
}

\bigskip

\normalsize

\noindent\begin{tabularx}{1\linewidth}{|p{11cm}C|}
  \hline
  \rule{0ex}{3ex}Établissement~:~\textbf{\etablissement} & Classe~:~\textbf{\classe}\\
  \rule[-1ex]{0ex}{4ex}NOM et Prénom du CANDIDAT~: \dotfill & Date de l'évaluation~:~\textbf{\datedeval}\\
  \hline
\end{tabularx}

\bigskip

\noindent{\Large Thème~: \theme}\\

\noindent{\Large Module(s) évalué(s)~:}

\begin{center}
  Capacités, connaissances et attitudes du référentiel évaluées\\
  \begin{tabularx}{16cm}{|C|l|}
    \hline
    \textbf{Capacités}     & \capacites\\
    \hline
    \textbf{Connaissances} & \connaissances\\
    \hline
    \textbf{Attitudes}     & \attitudes\\
    \hline
  \end{tabularx}
\end{center}

\noindent\fbox{\begin{minipage}[c]{.97\linewidth}
  \begin{description}
    \item[$\vartriangleright$] La clarté des raisonnements et la qualité de la rédaction
      interviendront dans l'appréciation des copies.
    \item[$\vartriangleright$] L'emploi des calculatrices est autorisé, dans les conditions
      prévues par la réglementation en vigueur.
  \end{description}
\end{minipage}}

\bigskip

L'évaluation est notée sur 10 : 7 points pour la mobilisation des
connaissances et des compétences pour résoudre des problèmes et 3
points pour l'activité d'expérimentation.

\begin{center}
  \begin{tabularx}{1\linewidth}{|L|c|}
    \hline
    \textbf{\rule{0ex}{3ex}Le candidat atteste avoir été informé de la date et des objectifs
      de l'évaluation le \makebox[4cm]{\rule{0ex}{3ex}\dotfill}} &
    \textbf{\rule[-6ex]{0ex}{3ex}\underline{Émargement}}\\
    \hline
  \end{tabularx}
\end{center}

%%%%%%%%%%%%%%%%%%%%%%%%%%%%%%%%%%%%%%%%%%%%%%%%%%%%%%%%%%
%                   fin page de garde
%%%%%%%%%%%%%%%%%%%%%%%%%%%%%%%%%%%%%%%%%%%%%%%%%%%%%%%%%%

\newpage

\lhead{\textbf{\matiere}}
\rhead{Année scolaire \annee}
\setcounter{page}{1}
\large

\noindent\begin{tabularx}{1\linewidth}[c]{|c|c|C|C|}
  \hline
  \multirow{2}{*}{\Large
  \begin{minipage}[c]{4cm}
    \begin{center}
      Contrôle en Cours\\ de Formation
    \end{center}
  \end{minipage}}
  &
  \multirow{2}{*}{
  \begin{minipage}[c]{7.2cm}
    \begin{center}
      Situation de Sciences Physiques et chimiques
    \end{center}
  \end{minipage}}
  &
  \rule[-1ex]{0ex}{4ex}Séquence &
  Durée\\
  \cline{3-4}
  &&
  \rule[-1.5ex]{0ex}{5ex}2/2
  &
  45~min\\
  \hline
\end{tabularx}
\medskip

\noindent\fbox{
  \Large\begin{minipage}[c]{.97\linewidth}
    \centerline{\rule[-.8ex]{0ex}{3.5ex}SUJET DESTINÉ AU CANDIDAT}
  \end{minipage}
}

\medskip

\noindent\begin{tabularx}{1\linewidth}{|p{10cm}L|}
  \hline
  \rule{0ex}{3ex}Établissement~:~\textbf{\etablissement} & Classe~:~\textbf{\classe}\\
  \rule[-1ex]{0ex}{4ex}NOM et Prénom~: \dotfill & Date de l'évaluation~: \textbf{\datedeval}\\
  \hline
\end{tabularx}

\medskip

\noindent\begin{tabularx}{1\linewidth}{|c|L|}
  \hline
  \multirow{2}{1.3cm}{\ \appel} &
  \textbf{\rule[-.8ex]{0ex}{3.5ex}L'examinateur intervient à la demande du
  candidat ou quand il le juge utile.}\\
  \cline{2-2}
  & \textbf{\rule[-.8ex]{0ex}{3.5ex}Dans la suite du document, ce symbole signifie
  \guillemotleft~Appeler l'examinateur~\guillemotright.}\\
  \hline
\end{tabularx}

\bigskip

\normalsize
\setstretch{1.2}
\parindent 0cm

\noindent
\begin{bclogo}[logo=\bcinfo]{Situation~:}
  \begin{center}
    \includegraphics[scale=0.4]{images/asterix.png}
  \end{center}
\end{bclogo}

\vfill

\begin{bclogo}{~Problématique~:}
  \begin{center}
    \textbf{Pourquoi voit-on toujours l'éclair avant d'entendre le tonnerre~?}
  \end{center}
\end{bclogo}

\vfill

\textbf{1. Compréhension et analyse de la situation / formulation d'une hypothèse.}\\[4pt]
\textbf{1.1. Proposer une explication à ce phénomène~:}\\[15pt]
\begin{minipage}[c]{.91\linewidth}
  .\dotfill\\

  .\dotfill\\
\end{minipage}\SA\A

\vfill

\textbf{1.2. Proposer une expérience permettant de mesurer la vitesse du son~:}\\[10pt]
\begin{minipage}[c]{.93\linewidth}
  \cadreblanc{40}
\end{minipage}\A\Ccomp   % CORRIGÉ : \C → \Ccomp

\newpage

\textbf{2. Expérimentation / modélisation de la situation}\\

\textbf{2.1. Réaliser le montage suivant}\\

\begin{center}
  \includegraphics[scale=0.6]{images/image1.png}
\end{center}  \Rcomp

\textbf{2.2 Mesures~:}
Compléter le tableau suivant~:
\renewcommand{\arraystretch}{1.8}
\begin{center}
  % CORRIGÉ : L{} → Lb{}, C{} → Cc{}
  \begin{tabular}{|Z{4cm}|Y{1cm}|Y{1cm}|Y{1cm}|Y{1cm}|Y{1cm}|Y{1cm}|Y{1cm}|Y{1cm}|Y{1cm}|}
    \hline
    Temps ($\mu$s) & & & & & & & & & \\   % CORRIGÉ : \micro → $\mu$
    \hline
    Distance (cm)  & 10 & 20 & 30 & 40 & 50 & 60 & 70 & 80 & 90\\
    \hline
  \end{tabular}
\end{center}

\vspace{3mm}

\textbf{2.3 Exploitation~:} \\[1mm]
\textbf{2.3.1 Tracé du graphique~:}\\[2mm]
Sur le repère ci-dessous, placer les points expérimentaux en prenant~:
\begin{itemize}
  \item en abscisse~: le temps $t$ (en $\mu$s)
  \item en ordonnée~: la distance $d$ (en cm)
\end{itemize}
Tracer ensuite la droite d'ajustement qui passe \textit{au mieux} par l'ensemble des points. \Rcomp

\vspace{3mm}

\begin{center}
  \begin{tikzpicture}
    \draw[very thin, gray!40] (0,0) grid[xstep=0.5, ystep=0.5] (12,9);
    \draw[thin, gray!70]      (0,0) grid[xstep=1,   ystep=1]   (12,9);
    % Axes
    \draw[->, thick] (0,0) -- (12.4,0) node[right] {$t$ ($\mu$s)};
    \draw[->, thick] (0,0) -- (0,9.4)  node[above] {$d$ (cm)};
    % Graduations x
    \foreach \x/\val in {0/0,1/500,2/1000,3/1500,4/2000,5/2500,
                          6/3000,7/3500,8/4000,9/4500,10/5000,11/5500,12/6000}{
      \draw (\x,0) -- (\x,-0.15) node[below, font=\small] {\val};
    }
    % Graduations y
    \foreach \y/\val in {0/0,1/10,2/20,3/30,4/40,5/50,6/60,7/70,8/80,9/90}{
      \draw (0,\y) -- (-0.15,\y) node[left, font=\small] {\val};
    }
    % Titres axes
    \node[below=8mm]          at (6,0)   {\textit{Temps} ($\mu$s)};
    \node[rotate=90,above=8mm] at (0,4.5) {\textit{Distance} (cm)};
  \end{tikzpicture}
\end{center}

\vspace{2mm}

\newpage

\textbf{2.3.2 Calcul du rapport $d/t$ pour chaque mesure~:}
\vspace{2mm}
\begin{itemize}
  \item La droite passe-t-elle par l'origine~? Justifier physiquement. \A

\cadreblanc{12}

\item Si la droite passe par l'origine, le rapport $\dfrac{d}{t}$ est constant et correspond à une vitesse.

Compléter le tableau ci-dessous~: (Arrondir le résultat à $0,0001$) \Rcomp

\vspace{3mm}
\begin{center}
  \renewcommand{\arraystretch}{2.5}
  \begin{tabular}{|Z{3cm}|Y{1cm}|Y{1cm}|Y{1cm}|Y{1cm}|Y{1cm}|Y{1cm}|Y{1cm}|Y{1cm}|Y{1cm}|}
    \hline
    Distance $d$ (cm) & 10 & 20 & 30 & 40 & 50 & 60 & 70 & 80 & 90 \\
    \hline
    Temps $t$ ($\mu$s) & & & & & & & & & \\
    \hline
    $\dfrac{d}{t}$ $\left(\dfrac{\text{cm}}{\mu\text{s}}\right)$ & & & & & & & & & \\
    \hline
  \end{tabular}
\end{center}

\vspace{5mm}

Calculer la \textbf{valeur moyenne} du rapport $\dfrac{d}{t}$~:

\[
  \overline{\left(\frac{d}{t}\right)}
  = \frac{1}{9}\Bigl(\ldots + \ldots + \ldots + \ldots + \ldots + \ldots + \ldots + \ldots + \ldots\Bigr)
  = \ldots\ldots\ldots.........\left(\frac{\text{cm}}{\mu\text{s}}\right)
\]

\end{itemize}

\vfill

 \textbf{2.3.3 Calcul de la vitesse du son dans l'air~:}\\[3mm]

Le coefficient directeur $a = \dfrac{d}{t}$ correspond à la vitesse du son $v_{\text{son}}$.\\[3mm]
Convertir votre résultat en m/s
(rappel~: $1\;\mu\text{s} = 10^{-6}\;\text{s}$ \quad et \quad $1\;\text{cm} = 10^{-2}\;\text{m}$)~:\\[3mm]
\textit{Attention~: le module mesure un aller-retour~! La distance réelle est la moitié de la distance affichée.}

\cadreblanc{22} \A

\vfill

\textbf{3.Conclusion}

\textbf{3.1. Cette valeur est-elle en accord avec celle déterminée en cours~?}\\[25pt]
\begin{minipage}[c]{.91\linewidth}
  .\dotfill\\
\end{minipage}\V

\vfill

\newpage

\textbf{3.2. En vous aidant du texte ci-dessous, déterminer, en Mach, la vitesse de la lumière \\
 ($v = 300\,000$ km/s).}\\

 \medskip

\begin{bclogo}[logo=\bcinfo,arrondi=.1,couleur=gray!10,barre=snake,ombre=true,epOmbre=0.15,couleurOmbre=black!30]{}
  Lorsqu'on dit qu'un avion passe le mur du son, c'est qu'il atteint et dépasse la vitesse du son dans l'air.
  L'appareil vole au moins à 340 m/s, équivalent à 1~224 km/h. On dit alors qu'il atteint Mach~1
  (les Mach indiquent la vitesse d'un corps par rapport à la vitesse du son~: Mach~1 = une fois cette
  vitesse~; Mach~2 = deux fois).
\end{bclogo}
\SA\Rcomp\Ccomp   % CORRIGÉ : \R → \Rcomp, \C → \Ccomp
\cadreblanc{25}

\bigskip   % CORRIGÉ : \vline hors tableau → \medskip

\begin{bclogo}[logo=\bcinfo,arrondi=.1,couleur=gray!10,barre=snake,ombre=true,epOmbre=0.15,couleurOmbre=black!30]{Où en sommes-nous~?}
  La sonde Parker en frôlant le soleil ce 29 avril 2022 à 532~000 km/h est devenue l'objet le plus rapide
  jamais conçu par l'humanité. Fabriquée par la NASA, elle peut parcourir plus de 140 kilomètres en une
  seconde. À cette vitesse, elle pourrait rejoindre la Lune depuis la Terre en seulement 43 minutes
  (source~: Futura Sciences).
\end{bclogo}

\medskip   % CORRIGÉ : \vline hors tableau → \medskip

\newpage

%%%%%%%%%%%%%%%%%%%%%%%%%%%%%%%%%%%%%%%%%%%%%%%%%%%%%%%%%%
%%              Fiche d'évaluation
%%%%%%%%%%%%%%%%%%%%%%%%%%%%%%%%%%%%%%%%%%%%%%%%%%%%%%%%%%

\small
\thispagestyle{empty}
\setstretch{1}

\noindent\begin{tabularx}{1\linewidth}[c]{|L|l|l|}
  \hline
  \multicolumn{3}{|c|}{\large
  \begin{minipage}[c]{.7\linewidth}
    \begin{center}
      \bfseries \rule{0ex}{3ex}GRILLE NATIONALE D'ÉVALUATION

      EN MATHÉMATIQUES ET EN

      SCIENCES PHYSIQUES ET CHIMIQUES\rule[-1ex]{0ex}{3ex}
    \end{center}
  \end{minipage}}\\
  \hline
  \rule[-1ex]{0ex}{4ex}Nom~: &
  \rule{.3cm}{0cm}Diplôme préparé~: \diplome\rule{.3cm}{0cm} &
  Séquence d'évaluation\footnote{Chaque séquence propose la résolution de
    problèmes issus du domaine professionnel ou de la vie courante. En
    mathématiques, elle comporte un ou deux exercices~; la résolution
    de l'un d'eux nécessite la mise en \oe{}uvre de capacités
    expérimentales.} \no\sequence \\
  \hline
\end{tabularx}\\[.3cm]

\textbf{1. Liste des capacités, connaissances et attitudes évaluées}
\begin{center}
  \begin{tabularx}{1\linewidth}{|C|l|}
    \hline
    \textbf{Capacités}     & \capacites\\
    \hline
    \textbf{Connaissances} & \connaissances\\
    \hline
    \textbf{Attitudes}     & \attitudes\\
    \hline
  \end{tabularx}
\end{center}

\medskip

\textbf{2. Évaluation\footnote{Des appels permettent de s'assurer de
    la compréhension du problème et d'évaluer le degré de maîtrise de
    capacités expérimentales et la communication orale. Il y en a au
    maximum 3 en sciences physiques et
    chimiques.\\
    {\bfseries En sciences physiques et chimiques~:} L'évaluation porte
    nécessairement sur des capacités expérimentales. 3 points sur
    10 sont consacrés aux questions faisant appel à la compétence
    \og{}Communiquer\fg{}.}}
\begin{center}
  \begin{tabularx}{.98\linewidth}[c]{|c|C|c|p{0.42cm}|p{0.42cm}|p{0.4cm}|c|}
    \hline
    \textbf{Compétences\footnote{L'ordre de présentation ne correspond
        pas à un ordre de mobilisation des compétences. La compétence
        \og{} Être autonome, Faire preuve d'initiative \fg{} est prise en
        compte au travers de l'ensemble des travaux réalisés. Les appels
        sont des moments privilégiés pour en apprécier le degré
        d'acquisition.}} &
    \textbf{Capacités} & \textbf{Questions} &
    \multicolumn{4}{c|}{
    \begin{minipage}[c]{2.8cm}
      \begin{center}
        \textbf{\rule{0ex}{2.2ex}Appréciation\\ du niveau\\ d'acquisition\rule[-1.2ex]{0ex}{2ex}}
      \end{center}
    \end{minipage}\footnote{Le professeur peut utiliser toute forme d'annotation lui permettant
    d'évaluer l'élève (le candidat) par compétences.}}\\
    \hline
    S'approprier &
    Rechercher, extraire et organiser l'information. &
    \begin{minipage}[c]{1.2cm}
      \begin{center}
        \rule{0ex}{2.2ex}1.1 ; 2.1 ; 2.3 ; 3.3\rule[-1.2ex]{0ex}{2ex}
      \end{center}
    \end{minipage}
    & & & & \\
    \hline
    Analyser &
    \begin{minipage}{1.0\linewidth}
      \begin{itemize}
        \item Émettre une conjecture, une hypothèse.
        \item Proposer une méthode de résolution, un protocole expérimental.
      \end{itemize}
    \end{minipage} &
    \begin{minipage}[c]{1.2cm}
      \begin{center}
        \rule{0ex}{2.2ex}1.1 ; 1.2 ; A2\rule[-1.2ex]{0ex}{2ex}
      \end{center}
    \end{minipage}
    & & & & \\
    \hline
    Réaliser &
    \begin{minipage}{1.0\linewidth}
      \begin{itemize}
        \item \rule{0ex}{2.2ex}Choisir une méthode de résolution, un protocole expérimental.
        \item Exécuter une méthode de résolution, expérimenter, simuler.\rule[-1.2ex]{0ex}{2ex}
      \end{itemize}
    \end{minipage} &
    \begin{minipage}[c]{1.2cm}
      \begin{center}
        \rule{0ex}{2.2ex}2.1 ; 2.3 ; 2.4 ; 3.1 ; 3.3\rule[-1.2ex]{0ex}{2ex}
      \end{center}
    \end{minipage}
    & & & & \\
    \hline
    Valider &
    \begin{minipage}{1.0\linewidth}
      \begin{itemize}
        \item \rule{0ex}{2.2ex}Contrôler la vraisemblance d'une conjecture, d'une hypothèse.
        \item Critiquer un résultat, argumenter.\rule[-1.2ex]{0ex}{2ex}
      \end{itemize}
    \end{minipage} &
    \begin{minipage}[c]{1.2cm}
      \begin{center}
        \rule{0ex}{2.2ex}2.2 ; 3.2\rule[-1.2ex]{0ex}{2ex}
      \end{center}
    \end{minipage}
    & & & & \\
    \hline
    Communiquer &
    \rule{0ex}{2.2ex}Rendre compte d'une démarche, d'un résultat, à l'oral ou à l'écrit.\rule[-1.2ex]{0ex}{2ex} &
    \begin{minipage}[c]{1.2cm}
      \begin{center}
        \rule{0ex}{2.2ex}1.2 ; A1 ; 3.2\rule[-1.2ex]{0ex}{2ex}
      \end{center}
    \end{minipage}
    & & & & \\
    \hline
    \multicolumn{3}{c|}{} & \multicolumn{4}{r|}{\rule[-1ex]{0ex}{4ex}\rule{3ex}{0ex}{\Large /~10}}\\
    \cline{4-7}
  \end{tabularx}
\end{center}

\end{document}